\section{Open questions}

The replication above confirms the paper's core mechanism qualitatively
and (in the small-instance regime) quantitatively. The following five
questions remain open and are candidates for a follow-up probe.

\begin{enumerate}

\item \textbf{Does the spoofer's advantage decay quantitatively as
  $15^{-d}$ predicted by Theorem~1.1, or does the effective base
  differ for the 1D brick-wall ensemble at small $n$?}

  \emph{Basis.} Our controls at $(n{=}6, d{=}1,2,3,6)$ show smooth
  decay of $F_{\mathrm{XEB}}$ from $\sim\!7.6\to 2.24\to 1.26\to 0.66$
  with a fixed 4-qubit block, and $(n{=}8)$ $13.88\to 2.31\to
  2.10\to 0.77$. The ratios do not obviously match a clean $1/15$ per
  depth step; a quantitative fit was not attempted.

  \emph{Next steps.} Sweep $d\in\{1,\ldots,8\}$ for
  $n\in\{8,10,12\}$ at matched block sizes $L\in\{2,4,6\}$, fit
  $\log F_{\mathrm{XEB}}$ vs $d$ and compare the slope to $\log 15$.
  Isolate the $(L/n)$ prefactor of Theorem~1.1 by adaptive light-cone
  block sizing.

\item \textbf{How does modern tensor-network / stabilizer-frame sampling
  change the spoofing threshold relative to the paper's light-cone
  algorithm?}

  \emph{Basis.} The paper (May 2020) predates the tensor-network
  breakthroughs (Pan-Zhang 2021, Kalachev et al.\ 2021, Liu et al.\
  2022) that classically sampled Sycamore-scale circuits without any
  light-cone restriction. Those methods may push the spoofable depth
  well beyond $\sqrt{\log n}$.

  \emph{Next steps.} Implement an MPS / MERA sampler (e.g.\ \texttt{quimb}
  or \texttt{opt\_einsum}) alongside the paper's spoofer on 1D
  brick-wall circuits up to $n{=}20$, $d{=}10$, and plot achievable
  $F_{\mathrm{XEB}}$ vs $(n,d)$ per strategy.

\item \textbf{At what physical two-qubit error rate does a noisy
  quantum sampler's $F_{\mathrm{XEB}}$ fall below what the light-cone
  spoofer can achieve on the same circuit?}

  \emph{Basis.} The paper is noiseless-classical: it shows
  $F_{\mathrm{XEB}}\gg 0$ is achievable classically without touching
  $q_C$. Sycamore's $F\approx 2.24\cdot 10^{-3}$ was consistent with
  per-gate errors $\sim 0.5\text{--}1\%$. We did not model noise.

  \emph{Next steps.} Add depolarising / amplitude-damping channels
  (density-matrix sim) with $p_2\in\{10^{-3}, 5\cdot 10^{-3}, 10^{-2},
  2\cdot 10^{-2}, 5\cdot 10^{-2}\}$, compare $F_{\mathrm{XEB}}$(noisy
  quantum) vs $F_{\mathrm{XEB}}$(light-cone spoofer) on the same
  circuit at each depth, and locate the crossover.

\item \textbf{Does a heavy-output-probability (HOP) benchmark resist
  shallow-circuit spoofing where Linear XEB does not?}

  \emph{Basis.} The attack exploits a specific structural feature of
  Linear XEB: any sampler that concentrates on high-$q_C$ bitstrings
  scores well. Alternative benchmarks (HOG, IQP-style heavy-output,
  Bouland-Fefferman-Nirkhe-Vazirani verified tests) may have different
  hardness profiles.

  \emph{Next steps.} Implement HOG on the same $n\le 8$ ensemble; for
  each circuit compute median $q_C$; measure fraction of samples with
  $q_C(x)>$ median. Compare exact / uniform / light-cone-spoofer HOG
  scores at $d=1,2,3,6$.

\item \textbf{How sensitive is the spoofability threshold to the gate
  set (Haar vs Sycamore fSim vs Clifford$+T$ vs iSWAP) and to circuit
  architecture (1D brick-wall vs 2D nearest-neighbour vs long-range
  random)?}

  \emph{Basis.} We tested only 1D brick-wall with Haar-random 2-qubit
  gates. Real supremacy experiments use structured gate sets (Sycamore
  fSim; IBM CX-basis with structured single-qubit layers) on 2D
  lattices, changing both light-cone growth and $t$-design convergence.

  \emph{Next steps.} Reimplement depth-1 and light-cone-4 spoofers on
  (a) 2D nearest-neighbour brick-wall at $n=9$ (3$\times$3), $n=16$
  (4$\times$4); (b) Sycamore-like fSim in place of Haar; (c)
  Clifford$+T$ (efficiently simulable in a different sense). Measure
  $F_{\mathrm{XEB}}$(spoofer) and Porter-Thomas convergence for each.

\end{enumerate}
